% =====================================================================
%  Chương 7 — Hệ đa tác tử (làn L6)
% =====================================================================
\chapter{Hệ đa tác tử (Multi-Agent Systems)}
\label{chap:mas}

Các chương trước xây dựng từng thành phần thông minh riêng lẻ: một mô hình phân loại, một mạng dự báo, một bộ truy xuất
tri thức, một mô hình ngôn ngữ trả lời có căn cứ. Hệ thống thực tế thường cần \emph{nhiều} thành phần như vậy phối hợp:
một thành phần lập kế hoạch, một thành phần tra cứu, một thành phần tính toán, một thành phần kiểm tra lại. Khi mỗi thành
phần có vai trò, quyền hạn và trạng thái riêng, và chúng trao đổi với nhau bằng thông điệp, ta có một \textbf{hệ đa tác tử}.
Chương này trình bày khái niệm tác tử, các mẫu tổ chức phối hợp, giao thức trao đổi, cách quản lý trạng thái và kiểm
chứng, các tiêu chí đánh giá, rồi trả lời bằng thực nghiệm một câu hỏi thực dụng: \emph{với một mô hình ngôn ngữ nhỏ chạy
cục bộ, tổ chức nhiều tác tử có thực sự giải toán tốt hơn một lời gọi duy nhất không, và phải trả giá bao nhiêu?}

% ---------------------------------------------------------------------
\section{Tác tử và sự phối hợp nhiều vai trò}
\label{sec:mas-agent}

Russell và Norvig định nghĩa \textbf{tác tử} (agent) là bất cứ thứ gì \emph{cảm nhận} môi trường qua cảm biến và
\emph{tác động} lên môi trường qua cơ cấu chấp hành \cite{ch7:russell2021}. Wooldridge bổ sung các tính chất của tác
tử thông minh: tự chủ (tự quyết hành động không cần con người can thiệp từng bước), phản ứng (đáp lại thay đổi của môi
trường), chủ động (theo đuổi mục tiêu) và có khả năng xã hội (tương tác với tác tử khác) \cite{ch7:wooldridge2009}.
Trong thế hệ hệ thống hiện nay, ``bộ não'' của tác tử thường là một mô hình ngôn ngữ lớn (LLM): nó đọc quan sát dưới dạng
văn bản, suy luận, và chọn hành động như gọi công cụ, gửi thông điệp hoặc trả lời. Vòng lặp
\emph{quan sát -- suy luận -- hành động -- ghi nhớ} (Hình~\ref{fig:ch7-loop}) lặp cho tới khi đạt điều kiện dừng.

\hinhtikz{ch7_agent_loop}{Vòng lặp của một tác tử dựa trên LLM: quan sát, suy luận, hành động trên môi trường và cập
nhật trạng thái}{ch7-loop}

Mô hình ReAct \cite{ch7:yao2023react} (đã dùng ở Tiểu luận 02) là dạng đơn giản nhất của vòng lặp này: LLM xen kẽ các
dòng \emph{Thought} (suy luận) và \emph{Action} (gọi công cụ), đọc \emph{Observation} rồi lặp lại. Một lời gọi LLM đơn
lẻ chưa phải là tác tử nếu không có cơ chế quản lý hành động và trạng thái; tương tự, viết nhiều tên vai trò trong cùng
một prompt chưa tạo ra nhiều tác tử độc lập. Một \textbf{hệ đa tác tử} (multi-agent system -- MAS) là hệ gồm nhiều tác tử
tương tác, có thể hợp tác, cạnh tranh hay thương lượng. Chương này tập trung vào hệ \emph{hợp tác} giải một nhiệm vụ
chung, trong đó mỗi tác tử phải được đặc tả rõ: nhận gì, được phép làm gì, lưu trạng thái nào, gửi gì cho ai
(Bảng~\ref{tab:ch7-roles}).

\begin{table}[H]\centering\caption{Các vai trò thường gặp trong hệ đa tác tử hợp tác}\label{tab:ch7-roles}
{\small\begin{tabular}{L{2.6cm}L{3.8cm}L{4.0cm}L{4.0cm}}
\toprule
Vai trò & Đầu vào & Sản phẩm & Quyền công cụ \\
\midrule
Planner & yêu cầu, ràng buộc & danh sách bước có điều kiện hoàn tất & không \\
Retriever & câu hỏi, thực thể & bằng chứng kèm ID nguồn & tìm kiếm, CSDL, KG \\
Solver / Executor & nhiệm vụ, kế hoạch & kết quả có cấu trúc & chạy mã, gọi API \\
Verifier / Critic & kết quả + bằng chứng & chấp nhận / lý do cần sửa & chạy lại phép kiểm \\
Coordinator & trạng thái toàn bộ & chuyển bước, biểu quyết, dừng & điều phối \\
\bottomrule
\end{tabular}}
\end{table}

Vì sao lại chia một nhiệm vụ cho nhiều tác tử? Có ba lý do có cơ sở:
\begin{enumerate}
  \item \textbf{Chuyên môn hoá}: mỗi vai trò có prompt, công cụ và bộ nhớ ngắn hơn, tập trung hơn; ngữ cảnh của một tác
        tử không bị lấp đầy bởi chi tiết của tác tử khác.
  \item \textbf{Đa dạng hoá lời giải}: các tác tử độc lập (khác prompt, khác phương pháp, khác nhiệt độ lấy mẫu) mắc lỗi
        ít tương quan; tổng hợp nhiều lời giải giảm phương sai -- cùng nguyên lý với tổ hợp mô hình ở Chương~2
        và self-consistency \cite{ch7:wang2023sc}.
  \item \textbf{Kiểm chứng chéo}: một tác tử kiểm tra kết quả của tác tử khác thay vì tự kiểm tra chính mình; nghiên
        cứu cho thấy LLM tự sửa lỗi suy luận của bản thân mà không có tín hiệu ngoài thường không hiệu quả
        \cite{ch7:huang2024selfcorrect}.
\end{enumerate}
Mặt trái cũng rõ: nhiều vai trò hơn nghĩa là nhiều lượt gọi LLM, nhiều token, độ trễ cao hơn, và thêm một loại lỗi mới
-- \emph{lỗi phối hợp} (hiểu sai thông điệp của nhau, vòng lặp không dứt, đồng thuận vào một đáp án sai). Phân tích
hơn một nghìn vết chạy của các khung đa tác tử phổ biến, Cemri và cộng sự \cite{ch7:cemri2025} thấy phần lớn thất bại đến
từ đặc tả vai trò và phối hợp giữa các tác tử chứ không phải từ năng lực của mô hình nền. Vì vậy quyết định dùng nhiều tác
tử phải được biện minh bằng \emph{kết quả nhiệm vụ trên mỗi đơn vị chi phí}, không bằng số lượng thông điệp hệ thống tạo ra.

% ---------------------------------------------------------------------
\section{Kiến trúc điều phối}
\label{sec:mas-patterns}

Cách các tác tử được nối với nhau quyết định luồng thông tin, điểm nghẽn và chế độ hỏng. Bốn mẫu tổ chức cơ bản
(Hình~\ref{fig:ch7-patterns}):

\hinhtikz{ch7_patterns}{Bốn mẫu tổ chức hệ đa tác tử: điều phối -- thợ, đường ống, tranh luận và bảng đen chung}{ch7-patterns}

\begin{itemize}
  \item \textbf{Điều phối -- thợ} (orchestrator--worker, hub-and-spoke): một Coordinator nhận yêu cầu, giao việc cho các
        thợ chuyên môn và tổng hợp kết quả. Dễ kiểm soát và ghi vết, nhưng Coordinator là điểm nghẽn và điểm hỏng đơn.
        AutoGen \cite{ch7:wu2023autogen} hiện thực mẫu này qua ``group chat manager''.
  \item \textbf{Đường ống} (pipeline / quy trình chuẩn): mỗi tác tử xử lý một công đoạn rồi chuyển tiếp, có thể có cạnh
        ngược để sửa (Verifier trả về Solver). MetaGPT \cite{ch7:hong2024metagpt} mã hoá quy trình phát triển phần mềm
        (product manager $\to$ architect $\to$ engineer $\to$ QA) theo cách này. Phù hợp khi quy trình đã biết trước.
  \item \textbf{Tranh luận} (debate): nhiều tác tử đưa lời giải, đọc lời giải của nhau và phản biện qua vài vòng, cuối
        cùng một giám khảo hoặc biểu quyết chốt đáp án. Du và cộng sự \cite{ch7:du2024debate} cho thấy tranh luận cải
        thiện độ đúng sự kiện và suy luận toán của LLM lớn.
  \item \textbf{Bảng đen chung} (blackboard): các tác tử không nhắn trực tiếp cho nhau mà đọc/ghi một kho trạng thái
        chung; một bộ điều khiển chọn tác tử nào được kích hoạt khi bảng thay đổi. Kiến trúc này có từ hệ nhận dạng tiếng
        nói Hearsay-II \cite{ch7:erman1980} và hợp với bài toán cần ghép nhiều mảnh tri thức không đồng bộ.
\end{itemize}

\begin{table}[H]\centering\caption{So sánh các mẫu điều phối}\label{tab:ch7-patterns}
{\small\begin{tabular}{L{2.9cm}L{3.6cm}L{4.0cm}L{4.0cm}}
\toprule
Mẫu & Luồng thông tin & Ưu điểm & Rủi ro chính \\
\midrule
Điều phối -- thợ & hình sao qua Coordinator & dễ ghi vết, dễ thêm thợ & điểm nghẽn, Coordinator hiểu sai \\
Đường ống & tuần tự, có vòng sửa & tái lập, dễ kiểm thử từng bước & lỗi sớm lan xuống cuối \\
Tranh luận & toàn đồ thị giữa các bên & đa dạng lời giải, tự phát hiện lỗi & tốn token, đồng thuận vào đáp án sai \\
Bảng đen & qua kho trạng thái & không đồng bộ, mở rộng tốt & xung đột ghi, khó suy luận về thứ tự \\
\bottomrule
\end{tabular}}
\end{table}

Hệ thực nghiệm ở mục~\ref{sec:mas-exp} kết hợp hai mẫu: \emph{đường ống} Planner $\to$ Solver $\to$ Verifier, trong đó
bước Solver gồm hai tác tử độc lập (một viết chương trình, một suy luận bằng lời), và khi chúng bất đồng thì chuyển sang
một vòng \emph{phân xử} kiểu tranh luận có biểu quyết.

% ---------------------------------------------------------------------
\section{Giao thức trao đổi}
\label{sec:mas-protocols}

Tác tử phối hợp được nhờ \emph{thông điệp}. Trong một tiến trình, thông điệp chỉ là một đối tượng (vai trò, nội dung,
siêu dữ liệu) nối vào lịch sử hội thoại. Khi tác tử chạy ở các tiến trình, máy chủ hoặc tổ chức khác nhau, cần giao thức
chuẩn. Hai giao thức mở đang được dùng rộng rãi giải quyết hai mối nối khác nhau (Hình~\ref{fig:ch7-protocols}):
\begin{itemize}
  \item \textbf{Model Context Protocol} (MCP) \cite{ch7:mcp2024} chuẩn hoá mối nối \emph{tác tử -- công cụ/dữ liệu}: một
        máy chủ MCP công bố danh sách công cụ (tên, mô tả, lược đồ JSON tham số), tài nguyên và mẫu prompt; tác tử (MCP
        client) liệt kê rồi gọi chúng qua JSON-RPC. Một công cụ viết một lần dùng được cho mọi tác tử hỗ trợ MCP.
  \item \textbf{Agent2Agent} (A2A) \cite{ch7:a2a2025} chuẩn hoá mối nối \emph{tác tử -- tác tử}: mỗi tác tử công bố một
        ``Agent Card'' mô tả kỹ năng và điểm truy cập; tác tử khác gửi \emph{task}, theo dõi trạng thái và nhận
        \emph{artifact} kết quả. A2A cho phép tác tử ``mờ'' (không lộ prompt, bộ nhớ, công cụ bên trong) vẫn hợp tác được.
\end{itemize}

\hinhtikz{ch7_protocols}{Hai giao thức bổ sung cho nhau: MCP nối tác tử với công cụ và dữ liệu, A2A nối tác tử với tác tử}{ch7-protocols}

Dù dùng giao thức nào, một thông điệp tốt giữa các tác tử nên có \emph{cấu trúc} (trường rõ ràng thay vì văn bản tự do),
mang \emph{nguồn gốc} (tác tử nào, lượt nào, dựa trên bằng chứng nào) và có \emph{điều kiện hoàn tất} kiểm tra được.
Thông điệp tự do giữa hai LLM là nguồn lỗi phối hợp lớn nhất: tác tử nhận có thể hiểu sai đơn vị, bỏ sót ràng buộc,
hoặc tin nhầm một con số do tác tử khác ``bịa''. Trong thực nghiệm, kết quả của Solver-PAL là một \emph{giá trị số} do
trình thông dịch Python trả về, không phải một câu văn do LLM viết -- đây là một cách đơn giản để thông điệp trở nên
kiểm chứng được.

% ---------------------------------------------------------------------
\section{Quản lý trạng thái và kiểm chứng}
\label{sec:mas-state}

Hệ đa tác tử phải trả lời ba câu hỏi về trạng thái: \emph{ai} giữ trạng thái (từng tác tử, Coordinator, hay một bảng đen
chung), \emph{cái gì} được lưu (lịch sử đầy đủ hay tóm tắt, kết quả trung gian, quyết định đã chốt), và \emph{bao lâu}
(một lượt, một phiên, hay bộ nhớ dài hạn). Các nguyên tắc thiết kế rút ra từ thực tiễn:
\begin{itemize}
  \item \textbf{Vết đầy đủ, có thứ tự}: mọi lượt gọi LLM và công cụ được ghi lại (vai trò, đầu vào, đầu ra, token, thời
        gian). Thiếu vết thì không thể đánh giá, gỡ lỗi hay quy trách nhiệm. Lớp \texttt{Trace} ở
        Đoạn mã~\ref{code:ch7_agents__Trace} làm việc này.
  \item \textbf{Cổng kiểm chứng tất định}: ở đâu kiểm tra được bằng chương trình (chạy lại phép tính, so khớp kiểu dữ
        liệu, kiểm tra đơn vị) thì dùng chương trình, không dùng LLM. LLM-Verifier chỉ dùng cho phần không kiểm tra tự
        động được.
  \item \textbf{Giới hạn vòng lặp và ngân sách}: số lượt sửa tối đa, số token tối đa, thời gian tối đa; khi vượt thì
        dừng và báo ``không tìm được đáp án'' thay vì trả một đáp án đoán.
  \item \textbf{Quyền tối thiểu cho công cụ}: tác tử chạy mã chỉ được làm đúng việc đó, trong hộp cát không mạng, không
        đọc tệp ngoài thư mục tạm, có giới hạn thời gian (mục~\ref{sec:mas-exp-design}).
\end{itemize}
Phương pháp Reflexion \cite{ch7:shinn2023} cho thấy một tác tử có thể cải thiện qua nhiều lượt khi nhận \emph{tín hiệu
ngoài} (kết quả kiểm thử, lỗi thực thi) và ghi bài học vào bộ nhớ; khi tín hiệu duy nhất là chính LLM tự nhận xét, mức
cải thiện giảm mạnh \cite{ch7:huang2024selfcorrect}. Thực nghiệm của chương áp dụng đúng nguyên lý này: lỗi thực thi
Python là tín hiệu ngoài được gửi lại cho Solver để sửa mã.

% ---------------------------------------------------------------------
\section{Đánh giá hệ đa tác tử}
\label{sec:mas-eval}

Độ đúng chỉ là một chiều. Một hệ đa tác tử chỉ đáng dùng khi phần cải thiện độ đúng bù được chi phí tăng thêm
(Hình~\ref{fig:ch7-eval}). Các chỉ số dùng trong thực nghiệm:
\begin{itemize}
  \item \textbf{Độ chính xác} (accuracy): tỷ lệ câu có đáp án số khớp đáp án vàng (sai số tuyệt đối $<10^{-4}$).
  \item \textbf{Số lượt gọi LLM} và \textbf{token} vào/ra mỗi câu: đại diện cho chi phí tính toán (với API trả phí, đây
        chính là tiền).
  \item \textbf{Độ trễ}: tổng thời gian sinh của LLM và thời gian chạy công cụ cho một câu.
  \item \textbf{Phân loại lỗi}: \emph{không có đáp án} (không trích được số, mã lỗi), \emph{lỗi số học} (một phép tính
        trong vết bị tính sai), \emph{lỗi suy luận} (hiểu sai đề, sai thứ tự bước), và riêng với hệ đa tác tử là
        \emph{lỗi phối hợp} -- có ít nhất một tác tử đã ra đáp án đúng nhưng cơ chế tổng hợp lại chọn sai.
  \item \textbf{Chi phí cho mỗi câu đúng thêm}: số lượt gọi (hoặc giây) phát sinh trên mỗi câu được sửa đúng so với hệ
        nền -- thước đo trực tiếp cho câu hỏi ``có đáng không''.
\end{itemize}

\hinhtikz{ch7_eval}{Các chiều đánh giá một hệ đa tác tử, từ độ đúng nhiệm vụ ở đáy đến an toàn ở đỉnh}{ch7-eval}

% ---------------------------------------------------------------------
\section{Rủi ro và giới hạn}
\label{sec:mas-risk}

\begin{table}[H]\centering\caption{Rủi ro của hệ đa tác tử và biện pháp giảm thiểu}\label{tab:ch7-risk}
{\small\begin{tabular}{L{3.6cm}L{5.0cm}L{5.8cm}}
\toprule
Rủi ro & Biểu hiện & Biện pháp \\
\midrule
Lỗi lan truyền & kế hoạch sai của Planner dẫn mọi Solver đi sai & Solver độc lập không phụ thuộc kế hoạch; Verifier đối chiếu đề gốc \\
Đồng thuận sai & nhiều tác tử cùng một LLM mắc cùng lỗi & đa dạng phương pháp (mã vs lời), kiểm chứng bằng công cụ \\
Vòng lặp vô hạn & tác tử sửa đi sửa lại không hội tụ & giới hạn số lượt, ngân sách token \\
Chi phí bùng nổ & mỗi câu tốn gấp nhiều lần & chỉ gọi tác tử phụ khi bất đồng (cổng đồng thuận) \\
Thực thi mã nguy hiểm & mã sinh ra đọc tệp, gọi mạng, treo máy & hộp cát: tiến trình riêng, chặn import, timeout \\
Tiêm lệnh (prompt injection) & dữ liệu ngoài chứa chỉ thị độc & tách dữ liệu khỏi chỉ thị, quyền tối thiểu \\
Mơ hồ trách nhiệm & không rõ tác tử nào gây lỗi & vết đầy đủ theo vai trò, ID lượt \\
\bottomrule
\end{tabular}}
\end{table}

Bộ phân loại lỗi MAST \cite{ch7:cemri2025}, xây dựng từ vết chạy của các khung đa tác tử phổ biến, cho thấy phần lớn
thất bại không đến từ năng lực của từng mô hình mà từ \emph{đặc tả} (vai trò chồng chéo, điều kiện dừng mơ hồ),
\emph{lệch hướng giữa các tác tử} (bỏ qua thông tin của nhau, tranh luận không hội tụ) và \emph{kiểm chứng yếu} (kết thúc
sớm, kiểm tra hời hợt). Ba nhóm này là danh sách kiểm tra khi thiết kế hệ ở mục tiếp theo: vai trò được viết rõ trong
lời nhắc hệ thống, điều kiện dừng là tất định (đồng thuận hoặc đã biểu quyết), và kiểm chứng dựa trên kết quả chạy mã.

% =====================================================================
\section{Thực nghiệm: đa tác tử giải toán lời văn GSM8K}
\label{sec:mas-exp}

\subsection{Bài toán, dữ liệu và câu hỏi nghiên cứu}

GSM8K \cite{ch7:cobbe2021} gồm 8.500 bài toán lời văn bậc tiểu học do con người viết; mỗi bài cần từ 2 đến 8 phép tính
cộng, trừ, nhân, chia, đáp án là một số. Bài toán ``dễ'' với con người nhưng là phép thử điển hình cho suy luận nhiều bước
của LLM: chỉ một phép tính sai hoặc một bước hiểu nhầm là sai cả câu. Thực nghiệm lấy một tập con cố định
\SL{ch7_gsm8k/n} câu của tập test (rút ngẫu nhiên với hạt giống 42, Đoạn mã~\ref{code:ch7_agents__load_gsm8k}); đáp án
vàng là con số sau ký hiệu \texttt{\#\#\#\#} trong lời giải mẫu. Một câu được chấm đúng khi số cuối cùng trích từ câu trả
lời khớp đáp án vàng với sai số tuyệt đối nhỏ hơn $10^{-4}$ (Đoạn mã~\ref{code:ch7_agents__to_number}).

\maNguon{ch7_agents__load_gsm8k}{Tập con cố định của GSM8K test và đáp án vàng}
\maNguon{ch7_agents__to_number}{Trích số cuối cùng trong câu trả lời (bỏ dấu phẩy ngăn nghìn, ký hiệu đô la)}

Mô hình ngôn ngữ duy nhất của mọi hệ là Qwen2.5-1.5B-Instruct \cite{ch7:qwen2024} -- cùng mô hình đã dùng ở Chương~5 --
chạy \emph{cục bộ} trên GPU RTX~4060 (bf16), giải mã tham lam (greedy) để mọi lần chạy cho cùng kết quả. Không dùng API
LLM bên ngoài. Chọn một mô hình nhỏ là có chủ ý: với mô hình lớn, GSM8K gần như đã bão hoà và khác biệt giữa các cách tổ
chức bị che khuất; với mô hình 1,5 tỷ tham số, mỗi câu trả lời đều ``mong manh'', nên tác dụng (hay tác hại) của việc
thêm tác tử hiện ra rõ. Ba câu hỏi nghiên cứu:
\begin{enumerate}
  \item[\textbf{Q1.}] Suy luận từng bước và dùng công cụ Python cải thiện độ đúng bao nhiêu so với trả lời trực tiếp?
  \item[\textbf{Q2.}] Hệ đa tác tử (Planner, hai Solver độc lập, Verifier, biểu quyết) có tốt hơn tác tử đơn tốt nhất
        không, và lỗi còn lại thuộc loại nào?
  \item[\textbf{Q3.}] Phần cải thiện (nếu có) phải trả bằng bao nhiêu lượt gọi, token và thời gian?
\end{enumerate}

\subsection{Thiết kế bốn hệ}
\label{sec:mas-exp-design}

Bốn hệ được so sánh trên đúng cùng tập câu hỏi, cùng LLM, cùng cách chấm (Hình~\ref{fig:ch7-systems}):
\begin{description}
  \item[Trực tiếp (direct)] một lượt gọi, lời nhắc hệ thống yêu cầu chỉ in ra con số -- mốc dưới, đo những gì mô hình
        ``nhớ'' hoặc đoán được mà không suy luận ra giấy.
  \item[CoT] một tác tử suy luận từng bước \cite{ch7:wei2022}, kết thúc bằng câu ``The answer is N.''.
  \item[PAL + Python] một tác tử viết hàm \texttt{solution()} bằng Python \cite{ch7:gao2023pal}; công cụ thực thi trong
        hộp cát trả về giá trị. Nếu mã lỗi, thông báo lỗi được gửi lại để tác tử sửa đúng một lần.
  \item[Đa tác tử (multi)] Planner lập kế hoạch; Solver-PAL viết mã theo kế hoạch, Solver-CoT giải bằng lời một cách độc
        lập; nếu hai đáp án khớp, Verifier chấp nhận ngay (cổng đồng thuận, 3 lượt gọi); nếu không, Verifier đọc cả hai
        lời giải để phân xử, một Solver-PAL thứ hai lấy mẫu ở nhiệt độ 0,7, rồi Coordinator biểu quyết đa số trên bốn
        đáp án (hoà phiếu thì tin Verifier) -- ý tưởng gần với self-consistency \cite{ch7:wang2023sc} và tranh luận đa
        tác tử \cite{ch7:du2024debate}.
\end{description}

\hinhtikz{ch7_systems}{Bốn hệ được so sánh, từ một lượt gọi đến một nhóm năm vai trò có cổng đồng thuận}{ch7-systems}

\hinhtikz{ch7_system_3d}{Kiến trúc hệ đa tác tử theo ba tầng: hạ tầng (GPU, LLM), tác tử và điều phối; bên phải là
công cụ hộp cát và dữ liệu}{ch7-system-3d}

Thiết kế này áp dụng các nguyên tắc ở mục~\ref{sec:mas-state}. Hai Solver dùng hai \emph{phương pháp khác nhau} (chương
trình và lời văn) chứ không chỉ hai lần gọi giống hệt nhau, để lỗi của chúng ít tương quan hơn -- hai lần gọi tham lam
cùng một lời nhắc sẽ cho cùng một đáp án và đồng thuận là vô nghĩa. Cổng đồng thuận là một kiểm tra \emph{tất định} (so
khớp hai số), nên phần lớn câu chỉ tốn 3 lượt gọi; vòng phân xử đắt hơn chỉ chạy khi thực sự có mâu thuẫn. Mọi thông
điệp giữa tác tử là văn bản kèm vai trò, được ghi vào một \texttt{Trace} chung (Đoạn mã~\ref{code:ch7_agents__Trace}).

\maNguon{ch7_agents__Trace}{Vết dùng chung: mọi bước được ghi kèm vai trò, cùng bộ đếm lượt gọi, token và thời gian}
\maNguon{ch7_agents__Agent}{Một tác tử là một vai trò (lời nhắc hệ thống) trên LLM dùng chung; mỗi lần \texttt{act} là
một lượt gọi được ghi vết}

Lời nhắc hệ thống của từng vai trò được giữ ngắn, viết bằng tiếng Anh (ngôn ngữ của đề bài), và mỗi lời nhắc nói rõ
\emph{định dạng đầu ra} mà bước sau cần: con số, câu ``The answer is N.'', hay một khối mã có hàm \texttt{solution()}.
Planner bị cấm tính đáp án cuối để không ``mớm'' một con số sai cho Solver.

\maNguon{ch7_agents__SYS_COT}{Lời nhắc hệ thống của tác tử CoT}
\maNguon{ch7_agents__SYS_PAL}{Lời nhắc hệ thống của tác tử PAL}
\maNguon{ch7_agents__SYS_PLANNER}{Lời nhắc của Planner: liệt kê đại lượng và các bước, không tính kết quả}
\maNguon{ch7_agents__SYS_CRITIC}{Lời nhắc của Verifier khi hai Solver bất đồng}
\maNguon{ch7_agents__PAL_SHOT}{Ví dụ mẫu (one-shot) gửi kèm mọi lượt gọi PAL}

\subsubsection{Tác tử PAL và công cụ thực thi trong hộp cát}

Hình~\ref{fig:ch7-pal} mô tả vòng PAL: LLM sinh mã, công cụ chạy mã, kết quả (hoặc lỗi) quay lại. Mã do LLM sinh là
\emph{mã không tin cậy}, nên công cụ được thiết kế theo quyền tối thiểu: mỗi lần chạy là một tiến trình Python con riêng
ở chế độ cô lập (\texttt{-I}, bỏ qua biến môi trường và thư mục người dùng), thư mục làm việc là một thư mục tạm bị xoá
sau khi chạy, biến môi trường gần như rỗng, giới hạn 5 giây; đoạn mở đầu (prelude) thay hàm \texttt{import} bằng bản chỉ
cho phép sáu thư viện tính toán và xoá \texttt{open}, \texttt{eval}, \texttt{exec}, \texttt{compile}, \texttt{input}. Mã
không thể đọc tệp, mở socket (\texttt{socket} không nằm trong danh sách) hay chạy mãi.

\hinhtikz{ch7_pal}{Vòng PAL: tác tử sinh chương trình, hộp cát thực thi, lỗi được phản hồi để tác tử tự sửa}{ch7-pal}

\maNguon{ch7_agents__SANDBOX_PRELUDE}{Đoạn mở đầu của hộp cát: danh sách trắng thư viện và xoá các hàm nguy hiểm}
\maNguon{ch7_agents__run_python}{Công cụ \texttt{python}: chạy mã không tin cậy trong tiến trình con cô lập, có thời hạn}
\maNguon{ch7_agents__extract_code}{Lấy khối mã Python đầu tiên trong câu trả lời của LLM}
\maNguon{ch7_agents__pal_once}{Tác tử PAL: sinh mã, gọi công cụ, gửi lỗi thực thi lại để tự sửa (tín hiệu ngoài)}

Điểm mấu chốt của \texttt{pal\_once} là phản hồi được gửi lại cho tác tử là \emph{thông báo lỗi thật} từ trình thông
dịch (ví dụ \texttt{NameError}, \texttt{ZeroDivisionError}, ``module not allowed''), không phải một nhận xét do LLM tự
viết -- đúng loại tín hiệu ngoài mà Reflexion cần \cite{ch7:shinn2023}.

\subsubsection{Ba hệ còn lại và cơ chế phối hợp}

\maNguon{ch7_agents__solve_direct}{Hệ trực tiếp: một lượt gọi, chỉ in con số}
\maNguon{ch7_agents__solve_cot}{Hệ CoT: một tác tử suy luận từng bước}
\maDai{ch7_agents__solve_multi}{Hệ đa tác tử: Planner, hai Solver độc lập, cổng đồng thuận, Verifier phân xử và
biểu quyết đa số}

Hình~\ref{fig:ch7-message-flow} là sơ đồ trình tự các thông điệp của một câu ở nhánh bất đồng. Nhánh đồng thuận dừng ngay
sau bước so khớp, nên chi phí trung bình của hệ phụ thuộc trực tiếp vào tỷ lệ câu mà hai Solver đồng ý với nhau.

\hinhtikz{ch7_message_flow}{Trình tự thông điệp của hệ đa tác tử ở nhánh bất đồng: kế hoạch, hai lời giải độc lập, phân
xử, mẫu bổ sung và biểu quyết}{ch7-message-flow}

\subsubsection{Chấm điểm và phân loại lỗi}

Mỗi câu của mỗi hệ được ghi thành một dòng JSON gồm đáp án, đúng/sai, số lượt gọi, token vào/ra, thời gian và toàn bộ
vết (Đoạn mã~\ref{code:ch7_agents__run_question}). Câu sai được phân loại tự động
(Đoạn mã~\ref{code:ch7_agents__classify_error}): \emph{không ra số} khi không trích được số hoặc mã không chạy;
\emph{sai số học} khi trong vết có một biểu thức dạng $a \circ b = c$ mà $c$ sai; \emph{lỗi phối hợp} (chỉ hệ đa tác tử)
khi có ít nhất một tác tử con đã đưa ra đúng đáp án vàng nhưng bước tổng hợp chọn khác; còn lại là \emph{sai lập luận}.
Khoảng tin cậy 95\,\% của độ chính xác được ước lượng bằng bootstrap trên các câu hỏi
(Đoạn mã~\ref{code:ch7_eval__bootstrap_ci}).

\maNguon{ch7_agents__classify_error}{Phân loại lỗi một câu sai}
\maNguon{ch7_agents__run_question}{Chạy một hệ trên một câu và ghi lại kết quả cùng vết}
\maNguon{ch7_eval__bootstrap_ci}{Khoảng tin cậy 95\,\% của accuracy bằng bootstrap}

\subsubsection{Chạy LLM cục bộ trên máy dùng chung}

Thực nghiệm chạy song song với các làn huấn luyện của những chương khác trên cùng một máy (một GPU 8\,GB, bộ nhớ ảo của
Windows có hạn). Hai sự cố thực tế đã định hình mã chạy. Thứ nhất, cách nạp mặc định của thư viện ánh xạ tệp
\texttt{safetensors} vào bộ nhớ; khi bộ nhớ commit của Windows gần cạn, tiến trình bị lỗi ``paging file too small'' hoặc
dừng đột ngột ngay lúc nạp. Bộ nạp tiết kiệm bộ nhớ dựng mô hình trên thiết bị \texttt{meta} (không cấp phát), đọc từng
tensor rồi chuyển thẳng lên GPU (Đoạn mã~\ref{code:ch7_llm__load_lowmem}). Thứ hai, GPU được chia sẻ: trên Windows, khi
VRAM không đủ, trình điều khiển có thể đặt trọng số sang bộ nhớ chia sẻ của hệ thống, và khi một tiến trình khác chạy
kernel liên tục, mỗi bước sinh token -- vốn gồm hàng trăm kernel nhỏ -- phải chờ lượt trên GPU. Cả hai đều không báo lỗi,
chỉ làm tốc độ sinh giảm mạnh, nên mã chờ đủ VRAM trống trước khi nạp (Đoạn mã~\ref{code:ch7_run__wait_for_vram}), ghi
từng câu ngay sau khi chấm để lần chạy sau tiếp tục từ chỗ dừng, và thời gian trong mục kết quả phải được đọc như số đo
\emph{tương đối} giữa các hệ trên cùng điều kiện máy, không phải độ trễ tuyệt đối của mô hình.

\maDai{ch7_llm__load_lowmem}{Nạp Qwen2.5-1.5B tiết kiệm bộ nhớ: dựng trên thiết bị meta rồi gán trọng số đọc từng tensor}
\maNguon{ch7_run__wait_for_vram}{Chờ đủ VRAM trống trước khi nạp, tránh trọng số rơi sang bộ nhớ chia sẻ}
\maNguon{ch7_run__run_system}{Chạy một hệ trên toàn bộ câu hỏi, ghi từng dòng kết quả và tiếp tục được khi bị ngắt}

% ---------------------------------------------------------------------
\section{Kết quả thực nghiệm}
\label{sec:mas-results}

\begin{luuy}{Cỡ mẫu}
LLM chạy trên một GPU dùng chung với các làn huấn luyện khác; tốc độ sinh thực tế chỉ khoảng
\SL{ch7_gsm8k/systems/cot/tok_per_s} token/giây, nên một câu của hệ đa tác tử mất nhiều phút, và lượt chạy phải nhường
GPU cho thực nghiệm của chương khác. Do đó số câu mỗi hệ đã chạy khác nhau: trả lời trực tiếp \SL{ch7_gsm8k/full/direct/n}
câu, CoT \SL{ch7_gsm8k/full/cot/n} câu, PAL \SL{ch7_gsm8k/full/pal/n} câu và hệ đa tác tử \SL{ch7_gsm8k/full/multi/n}
câu. Để so sánh công bằng, mọi bảng và hình dưới đây dùng \emph{tập chung} gồm \SL{ch7_gsm8k/n} câu mà cả bốn hệ đều đã
chạy -- mọi so sánh đều là so sánh cặp trên cùng một tập câu. Khoảng tin cậy bootstrap ở Bảng~\ref{tab:ch7-main} cho thấy
độ bất định tương ứng; các chênh lệch nhỏ hơn bề rộng khoảng này không nên được diễn giải như khác biệt thật.
\end{luuy}

\begin{table}[H]
\centering
\caption{Số câu đã chạy và độ chính xác của từng hệ trên toàn bộ số câu đó (tham khảo; so sánh chính dùng tập chung)}
\label{tab:ch7-full}
\begin{tabular}{lrr}
\toprule
Hệ & Số câu đã chạy & Accuracy (\%) \\
\midrule
Trả lời trực tiếp & \SL{ch7_gsm8k/full/direct/n} & \SL{ch7_gsm8k/full/direct/accPct} \\
CoT & \SL{ch7_gsm8k/full/cot/n} & \SL{ch7_gsm8k/full/cot/accPct} \\
PAL & \SL{ch7_gsm8k/full/pal/n} & \SL{ch7_gsm8k/full/pal/accPct} \\
Đa tác tử & \SL{ch7_gsm8k/full/multi/n} & \SL{ch7_gsm8k/full/multi/accPct} \\
\bottomrule
\end{tabular}
\end{table}

\subsection{Độ chính xác và chi phí}

\bang{ch7_main}{Độ chính xác, khoảng tin cậy 95\,\% và chi phí trung bình mỗi câu của bốn hệ trên GSM8K}{ch7-main}

\hinhpdf{ch7_acc}{1}{Độ chính xác của bốn hệ; thanh lỗi là khoảng tin cậy 95\,\% bootstrap}{ch7-acc}

\textbf{Q1 -- suy luận và công cụ.} Trả lời trực tiếp chỉ đúng \SL{ch7_gsm8k/systems/direct/accPct}\,\%: mô hình 1,5 tỷ
tham số không ``nhẩm'' được bài toán nhiều bước trong một con số. Cho phép suy luận từng bước nâng độ chính xác lên
\SL{ch7_gsm8k/systems/cot/accPct}\,\% (chênh lệch \SL{ch7_gsm8k/cot_vs_directPct} điểm) -- bước nhảy lớn nhất của toàn bộ
thực nghiệm, với chi phí chỉ là token sinh thêm trong \emph{một} lượt gọi. PAL đạt \SL{ch7_gsm8k/systems/pal/accPct}\,\%
(so với CoT: \SL{ch7_gsm8k/pal_vs_cotPct} điểm). Giao phép tính cho Python loại bỏ lỗi số học, nhưng với mô hình nhỏ, viết
một chương trình đúng định dạng là một kỹ năng riêng: ở lượt đầu, \SL{ch7_gsm8k/multi/first_nofunc_palPct}\,\% câu của
PAL không có hàm \texttt{solution()} và phải nhờ vòng tự sửa (Hình~\ref{fig:ch7-tool}).

\textbf{Q2 -- đa tác tử.} Hệ đa tác tử đạt \SL{ch7_gsm8k/systems/multi/accPct}\,\%, chênh lệch so với CoT là
\SL{ch7_gsm8k/vs_cot/gainPct} điểm và so với PAL là \SL{ch7_gsm8k/vs_pal/gainPct} điểm. Bảng~\ref{tab:ch7-vs} tách con số
ròng thành hai chiều: so với PAL, hệ đa tác tử sửa đúng \SL{ch7_gsm8k/vs_pal/fixed} câu PAL sai nhưng làm hỏng
\SL{ch7_gsm8k/vs_pal/broken} câu PAL đúng; so với CoT là \SL{ch7_gsm8k/vs_cot/fixed} câu sửa đúng và
\SL{ch7_gsm8k/vs_cot/broken} câu làm hỏng. Như vậy trên tập này, hệ đa tác tử \emph{không} vượt được tác tử CoT đơn lẻ:
nó chỉ hơn PAL -- Solver yếu hơn trong hai Solver của chính nó -- một khoảng nằm gọn trong khoảng tin cậy, và kém CoT vì
làm hỏng nhiều câu hơn số câu sửa được. Các mục sau truy nguyên vì sao: giá trị của kiểm chứng chéo là có thật ở nhánh
đồng thuận, nhưng nhánh bất đồng (\SL{ch7_gsm8k/multi/disagree_n} trên \SL{ch7_gsm8k/n} câu) được phân xử bởi chính LLM
nhỏ, và Solver-PAL trong hệ bị kế hoạch của Planner kéo lệch khỏi định dạng mã. Với cỡ mẫu này, chênh lệch giữa ba hệ có
suy luận đều nhỏ hơn bề rộng khoảng tin cậy; điều rút ra được chắc chắn là \emph{cơ chế} tác động, chưa phải độ lớn của nó.

\bang{ch7_vs}{Hệ đa tác tử so với hai tác tử đơn: số câu sửa đúng, làm hỏng và chi phí phụ trội}{ch7-vs}

\textbf{Q3 -- chi phí.} Mỗi câu, hệ đa tác tử dùng trung bình \SL{ch7_gsm8k/systems/multi/calls} lượt gọi LLM và
\SL{ch7_gsm8k/systems/multi/tokens} token, so với \SL{ch7_gsm8k/systems/cot/calls} lượt và
\SL{ch7_gsm8k/systems/cot/tokens} token của CoT; thời gian mỗi câu gấp \SL{ch7_gsm8k/vs_cot/seconds_ratio} lần CoT
(Hình~\ref{fig:ch7-cost}). Riêng con số thời gian cần đọc thận trọng: hệ đa tác tử chạy vào lúc máy dùng chung bớt tải
hơn (tốc độ sinh \SL{ch7_gsm8k/systems/multi/tok_per_s} so với \SL{ch7_gsm8k/systems/cot/tok_per_s} token/giây khi chạy
CoT), nên tỷ số thời gian đánh giá thấp chi phí thật; số lượt gọi và số token là thước đo chi phí đáng tin hơn.
Hình~\ref{fig:ch7-pareto} đặt bốn hệ lên mặt phẳng chi phí -- độ chính xác: một hệ chỉ đáng dùng nếu nằm trên biên
Pareto, tức là không có hệ nào vừa rẻ hơn vừa đúng hơn. Ở đây CoT nằm trên biên, còn hệ đa tác tử tốn nhiều token hơn mà
không đúng hơn.

\hinhpdf{ch7_cost}{1}{Chi phí trung bình mỗi câu: số lượt gọi LLM, số token (vào + ra) và thời gian}{ch7-cost}
\hinhpdf{ch7_pareto}{1}{Đánh đổi giữa chi phí (token mỗi câu, thang log) và độ chính xác; kích thước điểm tỷ lệ với số
lượt gọi}{ch7-pareto}
\hinhpdf{ch7_latency}{1}{Phân bố thời gian trả lời một câu của từng hệ (thang log)}{ch7-latency}
\hinhpdf{ch7_tokens}{1}{Phân bố số token sinh ra mỗi câu}{ch7-tokens}
\hinhpdf{ch7_calls}{1}{Phân bố số lượt gọi LLM mỗi câu}{ch7-calls}

Phân bố thời gian (Hình~\ref{fig:ch7-latency}) của hệ đa tác tử trải rộng nhất. Câu ở nhánh bất đồng cần trung bình
\SL{ch7_gsm8k/multi/disagree_calls} lượt gọi, so với \SL{ch7_gsm8k/multi/agree_calls} lượt ở nhánh đồng thuận (thêm
Verifier và một mẫu PAL nữa), và mất \SL{ch7_gsm8k/multi/disagree_seconds} giây so với \SL{ch7_gsm8k/multi/agree_seconds}
giây; chênh lệch thời gian nhỏ hơn chênh lệch lượt gọi vì tốc độ sinh của máy dùng chung dao động trong lúc chạy. Về số
lượt gọi (Hình~\ref{fig:ch7-calls}): mọi lượt vượt quá một của PAL đều đến từ vòng tự sửa mã; với hệ đa tác tử, phần lớn
lượt gọi đến từ chính thiết kế phối hợp (ba lượt khi đồng thuận, năm lượt khi bất đồng), còn vòng tự sửa của Solver-PAL
cộng thêm trung bình \SL{ch7_gsm8k/multi/repair_calls} lượt mỗi câu.

\subsection{Cổng đồng thuận}

\hinhpdf{ch7_gate}{1}{Cổng đồng thuận của hệ đa tác tử: tỷ lệ câu hai Solver khớp nhau (trái) và độ chính xác ở từng nhánh,
của cả hệ, và cận trên oracle (phải)}{ch7-gate}

Hai Solver độc lập đồng ý với nhau ở \SL{ch7_gsm8k/multi/agree_sharePct}\,\% số câu
(\SL{ch7_gsm8k/multi/agree_n} câu) và độ chính xác của nhánh này là \SL{ch7_gsm8k/multi/agree_accPct}\,\%
(Hình~\ref{fig:ch7-gate}): khi hai phương pháp khác nhau (chương trình và lời văn) cho cùng một con số, con số đó đáng tin
hơn hẳn mỗi phương pháp riêng lẻ -- đây là giá trị cốt lõi của kiểm chứng chéo. Ở \SL{ch7_gsm8k/multi/disagree_n} câu bất
đồng, Verifier và biểu quyết chỉ đạt \SL{ch7_gsm8k/multi/disagree_accPct}\,\%. Cận trên ``oracle'' -- tỷ lệ câu mà
\emph{ít nhất một} tác tử con đã có đáp án đúng -- là \SL{ch7_gsm8k/multi/oracle_accPct}\,\%; khoảng cách giữa con số này
và độ chính xác của cả hệ chính là phần mà cơ chế tổng hợp để lọt. Với một LLM 1,5 tỷ tham số làm Verifier, việc chọn
đúng giữa hai lời giải mâu thuẫn khó gần như việc tự giải bài toán, đúng như nhận định rằng LLM khó tự sửa lỗi suy luận
khi không có tín hiệu ngoài \cite{ch7:huang2024selfcorrect}.

Cổng đồng thuận cũng là thước đo \emph{độ tin cậy} miễn phí: hệ có thể báo kèm đáp án mức tin ``cao'' (hai Solver khớp)
hay ``thấp'' (phải phân xử). Trong ứng dụng thật, câu mức thấp nên được chuyển cho người kiểm tra -- một cách dùng đa tác
tử thường có lợi hơn cả việc cố nâng độ chính xác trung bình.

\subsection{Phân tích lỗi}

\bang{ch7_errors}{Phân loại câu sai của từng hệ}{ch7-errors}
\hinhpdf{ch7_errors}{1}{Cơ cấu lỗi của bốn hệ}{ch7-errors}

Cơ cấu lỗi (Bảng~\ref{tab:ch7-errors}, Hình~\ref{fig:ch7-errors}) khác nhau rõ giữa các hệ. Hệ trực tiếp sai gần như hoàn
toàn do lập luận -- nó không có bước nào để sai số học. CoT có \SL{ch7_gsm8k/systems/cot/errors/arithmetic} câu sai số học
(một phép tính viết sai kết quả trong lời giải), loại lỗi mà PAL loại bỏ được nhờ trình thông dịch; đổi lại, PAL có
\SL{ch7_gsm8k/systems/pal/errors/no_answer} câu \emph{không ra số} khi mã vẫn không chạy được sau lượt sửa, và phần lớn câu
sai của PAL vẫn là lỗi lập luận -- chương trình chạy đúng nhưng mô hình hoá sai bài toán. Hệ đa tác tử có \SL{ch7_gsm8k/systems/multi/errors/coordination}
câu \emph{lỗi phối hợp}: đáp án đúng đã xuất hiện ở một tác tử con nhưng bị biểu quyết hoặc Verifier gạt đi -- loại lỗi chỉ
tồn tại khi có nhiều tác tử.

\hinhpdf{ch7_tool}{1}{Kết quả các lần chạy mã trong hộp cát của PAL và hệ đa tác tử}{ch7-tool}

Hình~\ref{fig:ch7-tool} phân loại mọi lần gọi công cụ Python. Một phát hiện về phối hợp đáng chú ý: ở hệ đa tác tử, lượt
đầu của Solver-PAL thiếu hàm \texttt{solution()} ở \SL{ch7_gsm8k/multi/first_nofunc_multiPct}\,\% số câu, so với
\SL{ch7_gsm8k/multi/first_nofunc_palPct}\,\% của PAL đơn lẻ. Đọc vết cho thấy nguyên nhân: Planner, dù được dặn ``không
tính đáp án'', vẫn thường viết lời giải đầy đủ với tên bước riêng, và Solver-PAL đặt tên hàm theo kế hoạch (ví dụ
\texttt{calculate\_bolts}) thay vì theo quy ước \texttt{solution()} trong lời nhắc hệ thống. Thông tin thêm từ tác tử khác
đã \emph{lấn át} đặc tả định dạng -- đúng loại lỗi ``đặc tả và lệch hướng giữa tác tử'' mà MAST mô tả
\cite{ch7:cemri2025}. Vòng tự sửa bằng thông báo lỗi thật cứu được \SL{ch7_gsm8k/multi/rescued_multiPct}\,\% các trường
hợp này ở hệ đa tác tử (PAL đơn lẻ: \SL{ch7_gsm8k/multi/rescued_palPct}\,\%), nhưng mỗi lần cứu tốn thêm một lượt gọi.

\hinhpdf{ch7_only_correct}{0.8}{Số câu đúng ở hệ hàng nhưng sai ở hệ cột: mức độ bổ sung cho nhau giữa các hệ}{ch7-only}
\hinhpdf{ch7_difficulty}{1}{Độ chính xác theo độ khó (số phép tính trong lời giải mẫu)}{ch7-difficulty}

Ma trận ``đúng ở hệ này, sai ở hệ kia'' (Hình~\ref{fig:ch7-only}) cho thấy CoT và PAL bổ sung cho nhau: mỗi hệ đúng một
số câu mà hệ kia sai, nên kết hợp hai hệ (như hệ đa tác tử làm) có tiềm năng -- vấn đề nằm ở khâu chọn. Theo độ khó
(Hình~\ref{fig:ch7-difficulty}, Bảng~\ref{tab:ch7-difficulty}), trả lời trực tiếp và CoT giảm dần khi bài cần nhiều phép
tính hơn, còn PAL và hệ đa tác tử dao động ở các nhóm giữa rồi giảm mạnh ở nhóm khó nhất. Hai nhóm khó chỉ có vài câu mỗi
nhóm, nên đây là xu hướng định tính hơn là số đo chính xác.

\bang{ch7_difficulty}{Độ chính xác (\%) theo độ khó của bài toán}{ch7-difficulty}
\hinhpdf{ch7_running}{1}{Độ chính xác luỹ kế theo số câu đã chấm: con số ổn định dần khi mẫu lớn lên}{ch7-running}

\subsection{Vết suy luận mẫu}

Vết đầy đủ cho phép kiểm tra \emph{cách} hệ đi đến đáp án. Bảng~\ref{tab:ch7-answers} đặt đáp án của bốn hệ cạnh nhau cho
vài câu mà các hệ cho kết quả khác nhau; các hộp vết dưới đây được chọn tự động từ dữ liệu chạy.

Hộp vết thứ nhất là một câu ở nhánh đồng thuận: Solver-PAL và Solver-CoT cho cùng một số nên hệ chốt ngay, không cần
phân xử. Hộp thứ hai là câu bất đồng mà PAL đơn lẻ trả lời sai nhưng hệ đa tác tử vẫn chọn đúng sau phân xử. Hộp thứ ba
là một lỗi phối hợp điển hình: Solver-PAL (sau một lần tự sửa) đã tính đúng, nhưng Solver-CoT và Verifier cùng ra một số
sai, mẫu PAL thứ hai lại thiếu hàm \texttt{solution()}, nên biểu quyết đa số chọn đáp án sai -- hai phiếu sai tương quan
thắng một phiếu đúng. Hộp cuối là PAL đơn lẻ tự sửa: lần chạy đầu báo lỗi, thông báo lỗi được đưa lại cho mô hình và lần
chạy sau cho kết quả đúng.

\begin{vetsuyluan}{Vết tác tử -- hai Solver đồng thuận, Verifier chấp nhận ngay (câu 1, đáp án 3, hệ trả lời 3, 4 lượt gọi, 208{,}5 s)}
\vetBuoc{cNavy}{CÂU HỎI}{A robe takes 2 bolts of blue fiber and half that much white fiber.  How many bolts in total does it take?}
\vetBuoc{cIndigo}{PLANNER}{1. Determine the amount of blue fiber: - A robe requires 2 bolts of blue fiber. 2. Determine the amount of white fiber: - The amount of white fiber is half the amount of blue fiber. - Therefore, the amount of white fiber is (2)/(2) = 1 bolt. 3. Calculate the total number of bolts: - Total bolts = Blue fiber + White fiber - Total bolts = 2 + 1 = 3 4. Final answer: - The total number of bolts required is 3.}
\vetBuoc{cGreen}{SOLVER-PAL}{def calculate\_bolts(): blue\_fiber = 2 white\_fiber = blue\_fiber / 2 total\_bolts = blue\_fiber + white\_fiber return total\_bolts}
\vetBuoc{cOrange}{PYTHON}{lỗi: không có hàm solution()}
\vetBuoc{cGreen}{SOLVER-PAL}{\texttt{\footnotesize def~solution():}\\
\texttt{\footnotesize ~~~~blue\_fiber~=~2}\\
\texttt{\footnotesize ~~~~white\_fiber~=~blue\_fiber~/~2}\\
\texttt{\footnotesize ~~~~total\_bolts~=~blue\_fiber~+~white\_fiber}\\
\texttt{\footnotesize ~~~~return~total\_bolts}}
\vetBuoc{cOrange}{PYTHON}{kết quả = 3}
\vetBuoc{cBlue}{SOLVER-CoT}{To solve this problem, we need to calculate the total number of bolts of fiber required for the robe. 1. Determine the amount of blue fiber: - The robe requires 2 bolts of blue fiber. 2. Determine the amount of white fiber: - The robe requires half as much white fiber as blue fiber. - Therefore, the amount of white fiber needed is (1)/(2) × 2 = 1 bolt. 3. Calculate the total number of bolts: - Add the number of blue fiber bolts to the number of white fiber bolts: 2 (blue) + 1 (white) = 3 bolts 4. State the final answer: - The total number of bolts of fiber required is 3. The answer is 3.}
\vetBuoc{cRed}{VERIFIER}{hai solver đồng thuận: 3 -> chấp nhận}
\end{vetsuyluan}


\bang{ch7_answers}{Đáp án của bốn hệ cho một số câu mà các hệ bất đồng}{ch7-answers}

\begin{vetsuyluan}{Vết tác tử -- hai Solver bất đồng, Verifier và biểu quyết chọn đúng (PAL đơn lẻ sai) (câu 43, đáp án 48, hệ trả lời 48, 6 lượt gọi, 144{,}9 s)}
\vetBuoc{cNavy}{CÂU HỎI}{According to its nutritional info, a bag of chips has 250 calories per serving. If a 300g bag has 5 servings, how many grams can you eat if your daily calorie target is 2000 and you have already consumed 1800 calories?}
\vetBuoc{cIndigo}{PLANNER}{1. Total calories in a bag: 250 calories/serving * 5 servings = 1250 calories 2. Daily calorie target: 2000 calories 3. Calories already consumed: 1800 calories 4. Remaining calories for today: 2000 calories - 1800 calories = 200 calories 5. Calories per gram of chips: 1250 calories / 300 grams = 4.17 grams/calorie 6. Remaining grams of chips: 200 calories / 4.17 grams/calorie ≈ 48.1 grams}
\vetBuoc{cGreen}{SOLVER-PAL}{def calculate\_remaining\_chips(): total\_calories\_per\_bag = 250 * 5 daily\_calorie\_target = 2000 calories\_consumed = 1800 remaining\_calories = daily\_calorie\_target - calories\_consumed calories\_per\_gram = total\_calories\_per\_bag / 300 remaining\_grams = remaining\_calories / calories\_per\_gram return round(remaining\_grams)}
\vetBuoc{cOrange}{PYTHON}{lỗi: không có hàm solution()}
\vetBuoc{cGreen}{SOLVER-PAL}{\texttt{\footnotesize def~solution():}\\
\texttt{\footnotesize ~~~~total\_calories\_per\_bag~=~250~*~5}\\
\texttt{\footnotesize ~~~~daily\_calorie\_target~=~2000}\\
\texttt{\footnotesize ~~~~calories\_consumed~=~1800}\\
\texttt{\footnotesize ~~~~remaining\_calories~=~daily\_calorie\_target~-~calories\_consumed}\\
\texttt{\footnotesize ~~~~calories\_per\_gram~=~total\_calories\_per\_bag~/~300}\\
\texttt{\footnotesize ~~~~remaining\_grams~=~remaining\_calories~/~calories\_per\_gram}\\
\texttt{\footnotesize ~~~~return~round(remaining\_grams)}}
\vetBuoc{cOrange}{PYTHON}{kết quả = 48}
\vetBuoc{cBlue}{SOLVER-CoT}{First, calculate the total calories in the bag: Total calories = Calories per serving × Number of servings Total calories = 250 × 5 = 1250 calories Next, determine how many more calories you need to reach your daily target: Remaining calories needed = Daily calorie target - Calories already consumed Remaining calories needed = 2000 - 1800 = 200 calories Now, find out how many servings correspond to these remaining calories: Servings needed = (Remaining calories needed)/(Calories per serving) Servings needed = (200)/(250) = 0.8 servings Finally, calculate the weight of the additional servings:  ... [lược bớt]}
\vetBuoc{cRed}{VERIFIER}{The answer is N.}
\vetBuoc{cTeal}{PAL\#2}{def calculate\_remaining\_calories(): total\_calories\_in\_bag = 250 * 5 daily\_calorie\_target = 2000 calories\_consumed\_today = 1800 remaining\_calories\_for\_today = daily\_calorie\_target - calories\_consumed\_today calories\_per\_gram = total\_calories\_in\_bag / 300 remaining\_grams\_of\_chips = remaining\_calories\_for\_today / calories\_per\_gram return remaining\_grams\_of\_chips Calculate the remaining grams of chips remaining\_grams\_of\_chips = calculate\_remaining\_calories() print(remaining\_grams\_of\_chips)}
\vetBuoc{cOrange}{PYTHON}{lỗi: không có hàm solution()}
\vetBuoc{cPurple}{ĐIỀU PHỐI}{phiếu \{48: 1, 1800: 1\} -> chọn 48}
\end{vetsuyluan}

\begin{vetsuyluan}{Vết tác tử -- lỗi phối hợp -- một tác tử đã ra đáp án đúng nhưng bị biểu quyết loại (câu 39, đáp án 18, hệ trả lời 16, 6 lượt gọi, 241{,}5 s)}
\vetBuoc{cNavy}{CÂU HỎI}{Dana can run at a rate of speed four times faster than she can walk, but she can skip at a rate of speed that is half as fast as she can run. If she can skip at 3 miles per hour, how many miles can she travel in six hours if she spends one-third of the time running and two-thirds of the time walking?}
\vetBuoc{cIndigo}{PLANNER}{1. Determine Dana's running speed: - Given that skipping at 3 miles per hour is half as fast as running. - Therefore, running speed = 2 × 3 = 6 miles per hour. 2. Determine Dana's walking speed: - Given that running at 6 miles per hour is four times faster than walking. - Therefore, walking speed = 6 / 4 = 1.5 miles per hour. 3. Calculate total distance for each activity: - Time spent running = 6 × (1)/(3) = 2 hours. - Distance running = 2 × 6 = 12 miles. - Time spent walking =}
\vetBuoc{cGreen}{SOLVER-PAL}{Given data skip\_speed = 3 \# miles per hour walk\_speed = skip\_speed / 4 \# miles per hour run\_speed = skip\_speed * 4 \# miles per hour Total time available total\_time = 6 \# hours Time spent on each activity time\_running = total\_time / 3 \# hours time\_walking = total\_time - time\_running \# hours Calculate distances distance\_running = time\_running * run\_speed \# miles distance\_walking = time\_walking * walk\_speed \# miles Total distance traveled total\_distance = distance\_running + distance\_walking \# miles}
\vetBuoc{cOrange}{PYTHON}{lỗi: không có hàm solution()}
\vetBuoc{cGreen}{SOLVER-PAL}{\texttt{\footnotesize It~seems~there~was~an~error~in~the~previous~code~because~it~did~not~define}\\
\texttt{\footnotesize ~~~~~the~`solution`~function~correctly.~Let~me~correct~this~and~provide~}\\
\texttt{\footnotesize ~~~~the~proper~solution:}\\
\texttt{\footnotesize def~solution():}\\
\texttt{\footnotesize ~~~~\#~Given~data}\\
\texttt{\footnotesize ~~~~total\_time~=~6~~\#~Total~time~in~hours}\\
\texttt{\footnotesize ~~~~time\_running\_fraction~=~1/3~~\#~Fraction~of~time~spent~running}\\
\texttt{\footnotesize ~~~~time\_walking\_fraction~=~2/3~~\#~Fraction~of~time~spent~walking}\\
\texttt{\footnotesize ~~~~\#~Calculate~time~spent~on~each~activity}\\
\texttt{\footnotesize ~~~~time\_running~=~total\_time~*~time\_running\_fraction}\\
\texttt{\footnotesize ~~~~time\_walking~=~total\_time~*~time\_walking\_fraction}\\
\texttt{\footnotesize ~~~~\#~Given~speeds}\\
\texttt{\footnotesize ~~~~running\_speed~=~6~~\#~Running~speed~in~miles~per~hour}\\
\texttt{\footnotesize ~~~~walking\_speed~=~running\_speed~/~4~~\#~Walking~speed~in~miles~per~hour}\\
\texttt{\footnotesize ~~~~\#~Calculate~distances}\\
\texttt{\footnotesize ~~~~distance\_running~=~time\_running~*~running\_speed}\\
\texttt{\footnotesize ~~~~distance\_walking~=~time\_walking~*~walking\_speed}\\
\texttt{\footnotesize ~~~~\#~Calculate~total~distance~traveled}\\
\texttt{\footnotesize ~~~~total\_distance~=~distance\_running~+~distance\_walking}\\
\texttt{\footnotesize ~~~~return~total\_distance}\\
\texttt{\footnotesize \#~Call~the~function~to~get~the~result}\\
\texttt{\footnotesize result~=~solution()}\\
\texttt{\footnotesize print(result)}\\
\texttt{\footnotesize This~code~will~calculate~the~total~distance~Dana~travels~in~six~hours,~}\\
\texttt{\footnotesize ~~~~given~the~fractions~of~time~spent~running~and~walking,~and~the~speeds~}\\
\texttt{\footnotesize ~~~~for~both~activities.}}
\vetBuoc{cOrange}{PYTHON}{kết quả = 18}
\vetBuoc{cBlue}{SOLVER-CoT}{To solve this problem, we need to determine Dana's speeds for running, walking, and skipping, then calculate the distance she travels in six hours based on her spending different portions of that time on each activity. 1. Determine Dana's running speed: - Skip rate = 3 mph (half as fast as running) - Running rate = 2 * Skip rate = 2 * 3 mph = 6 mph 2. Calculate the time spent on each activity: - Total time = 6 hours - Time spent running = 1/3 of total time = 6 hours / 3 = 2 hours - Time spent walking = 2/3 of total time = 6 hours * 2/3 = 4 hours 3. Calculate the distance traveled during each a ... [lược bớt]}
\vetBuoc{cRed}{VERIFIER}{The answer is 16.}
\vetBuoc{cTeal}{PAL\#2}{Given data walking\_speed = 1.5 \# miles per hour running\_speed = 6 \# miles per hour skip\_speed = 3 \# miles per hour Calculate the total time spent on each activity time\_running = 6 * (1/3) time\_walking = 6 * (2/3) Calculate the distance traveled by each activity distance\_running = time\_running * running\_speed distance\_walking = time\_walking * walking\_speed distance\_skipping = time\_running * skip\_speed Total distance traveled total\_distance = distance\_running + distance\_walking + distance\_skipping}
\vetBuoc{cOrange}{PYTHON}{lỗi: không có hàm solution()}
\vetBuoc{cPurple}{ĐIỀU PHỐI}{phiếu \{18: 1, 16: 2\} -> chọn 16}
\end{vetsuyluan}

\begin{vetsuyluan}{Vết tác tử -- PAL tự sửa sau khi mã chạy lỗi trong sandbox (câu 30, đáp án 109, hệ trả lời 109, 2 lượt gọi, 119{,}8 s)}
\vetBuoc{cNavy}{CÂU HỎI}{Darrell and Allen's ages are in the ratio of 7:11. If their total age now is 162, calculate Allen's age 10 years from now.}
\vetBuoc{cGreen}{PAL}{def calculate\_allens\_age(): total\_age = 162 ratio\_darrell = 7 ratio\_allen = 11 combined\_ratio = ratio\_darrell + ratio\_allen allen\_age = (ratio\_allen / combined\_ratio) * total\_age return int(allen\_age)}
\vetBuoc{cOrange}{PYTHON}{lỗi: không có hàm solution()}
\vetBuoc{cGreen}{PAL}{\texttt{\footnotesize def~solution():}\\
\texttt{\footnotesize ~~~~\#~Given~ratio~of~Darrell's~to~Allen's~age}\\
\texttt{\footnotesize ~~~~ratio\_darrell\_allen~=~[7,~11]}\\
\texttt{\footnotesize ~~~~\#~Total~age~of~Darrell~and~Allen}\\
\texttt{\footnotesize ~~~~total\_age~=~162}\\
\texttt{\footnotesize ~~~~\#~Calculate~the~sum~of~the~ratio~parts}\\
\texttt{\footnotesize ~~~~sum\_ratio\_parts~=~sum(ratio\_darrell\_allen)}\\
\texttt{\footnotesize ~~~~\#~Calculate~Darrell's~age}\\
\texttt{\footnotesize ~~~~darrell\_age~=~(ratio\_darrell\_allen[0]~/~sum\_ratio\_parts)~*~total\_age}\\
\texttt{\footnotesize ~~~~\#~Calculate~Allen's~age}\\
\texttt{\footnotesize ~~~~allen\_age~=~total\_age~-~darrell\_age}\\
\texttt{\footnotesize ~~~~\#~Calculate~Allen's~age~10~years~from~now}\\
\texttt{\footnotesize ~~~~allen\_age\_10\_years\_later~=~allen\_age~+~10}\\
\texttt{\footnotesize ~~~~return~allen\_age\_10\_years\_later}\\
\texttt{\footnotesize \#~Call~the~function~and~print~the~result}\\
\texttt{\footnotesize print(solution())}}
\vetBuoc{cOrange}{PYTHON}{kết quả = 109}
\end{vetsuyluan}


% ---------------------------------------------------------------------
\section{Triển khai hệ đa tác tử lên Web App}
\label{sec:mas-deploy}

Bốn hệ được đưa vào Web App chung của tiểu luận qua mô-đun \texttt{tl3/serve\_ch7.py}. \texttt{SPECS} khai báo một trải
nghiệm kiểu ``chat'': ô nhập bài toán (mặc định là bài đầu tiên của GSM8K test) và ô chọn hệ giải (\texttt{multi},
\texttt{pal}, \texttt{cot}, \texttt{direct}). Hàm \texttt{predict} gọi đúng hàm \texttt{solve\_*} đã dùng trong thực
nghiệm, rồi trả về đáp án, toàn bộ vết theo vai trò (mỗi bước là một cặp \emph{vai trò -- nội dung}, gồm cả kết quả công
cụ Python), danh sách công cụ và số lượt gọi, token. LLM được nạp lười bằng bộ nạp tiết kiệm bộ nhớ ở lần gọi đầu và giữ
trong bộ nhớ cho các lần sau. Giao diện hiển thị vết theo đúng thứ tự Planner $\to$ Solver $\to$ Python $\to$ Verifier
$\to$ Coordinator, nên người dùng thấy được \emph{vì sao} hệ chọn đáp án, không chỉ đáp án.

\maNguon{serve_ch7__SPECS}{Khai báo trải nghiệm \texttt{ch7\_math} cho Web App}
\maNguon{serve_ch7__predict}{Hàm suy luận: bài toán $\to$ đáp án, vết từng tác tử, công cụ đã dùng}

\anhUI{ch7_math_web}{Hệ đa tác tử giải bài toán lời văn trên Web App (desktop): đáp án, vết từng vai trò và công cụ}{ch7-ui-web}
\anhMobile{ch7_math_mobile}{Hệ đa tác tử trên khung điện thoại 390 px: đáp án và vết suy luận xếp dọc}{ch7-ui-mobile}

\begin{luuy}{Giới hạn sử dụng}
Mã do LLM sinh luôn chạy trong hộp cát (tiến trình riêng, không mạng, không đọc tệp, giới hạn thời gian). Mô hình 1,5 tỷ
tham số vẫn sai một phần đáng kể các bài toán; vết suy luận giúp kiểm tra lại, không thay cho việc kiểm tra.
\end{luuy}

% ---------------------------------------------------------------------
\section{Kết luận chương}

Chương này đi từ khái niệm tác tử, các mẫu kiến trúc và giao thức phối hợp tới một thực nghiệm có kiểm soát: bốn hệ dùng
chung một LLM cục bộ 1,5 tỷ tham số và cùng tập câu GSM8K, chỉ khác nhau ở cách tổ chức suy luận. Trên
\SL{ch7_gsm8k/n} câu chung:
\begin{itemize}
\item \textbf{Suy luận từng bước là đòn bẩy lớn nhất:} CoT nâng độ chính xác từ
\SL{ch7_gsm8k/systems/direct/accPct}\,\% lên \SL{ch7_gsm8k/systems/cot/accPct}\,\% chỉ với một lượt gọi.
\item \textbf{Công cụ không tự động tốt hơn:} PAL loại bỏ lỗi số học nhưng chỉ đạt \SL{ch7_gsm8k/systems/pal/accPct}\,\%,
vì mô hình nhỏ phải viết đúng cả chương trình lẫn định dạng.
\item \textbf{Đa tác tử không vượt tác tử đơn tốt nhất:} hệ đa tác tử đạt \SL{ch7_gsm8k/systems/multi/accPct}\,\% với
\SL{ch7_gsm8k/systems/multi/calls} lượt gọi và \SL{ch7_gsm8k/systems/multi/tokens} token mỗi câu; nó hơn PAL nhưng kém CoT,
và mọi chênh lệch đều nằm trong khoảng tin cậy.
\item \textbf{Giá trị thật của đa tác tử là tín hiệu tin cậy:} khi hai Solver khớp nhau
(\SL{ch7_gsm8k/multi/agree_sharePct}\,\% số câu), độ chính xác đạt \SL{ch7_gsm8k/multi/agree_accPct}\,\%; nhánh bất đồng
chỉ đạt \SL{ch7_gsm8k/multi/disagree_accPct}\,\% dù cận trên oracle là \SL{ch7_gsm8k/multi/oracle_accPct}\,\%.
\item \textbf{Phối hợp sinh ra lỗi mới:} kế hoạch của Planner lấn át đặc tả định dạng của Solver-PAL
(\SL{ch7_gsm8k/multi/first_nofunc_multiPct}\,\% lượt chạy đầu thiếu \texttt{solution()}), và
\SL{ch7_gsm8k/systems/multi/errors/coordination} câu sai vì khâu tổng hợp gạt đi đáp án đúng đã có trong vết.
\end{itemize}
Bài học thiết kế: với mô hình nhỏ, hãy bắt đầu từ tác tử đơn mạnh nhất (ở đây là CoT); chỉ thêm tác tử khi có một tín hiệu
kiểm chứng \emph{bên ngoài} đáng tin (trình thông dịch, bộ kiểm thử, người duyệt); và dùng sự bất đồng giữa các tác tử để
chuyển câu khó cho người kiểm tra thay vì để chính LLM yếu tự phân xử. Giới hạn của chương: cỡ mẫu nhỏ do tài nguyên tính
toán dùng chung, và chỉ thử một LLM. Dùng một mô hình lớn hơn làm Verifier, hoặc cho Verifier chạy lại cả hai chương trình
thay vì đọc lời văn, là hướng mở rộng tự nhiên.
